% =============================================================================
% CHAPTER 2 — THEORETICAL FOUNDATION (REDUCED, six-section version)
% Thesis: PATHCAST
% Style:  formal definition -> central property -> deferral to appendix/chapter
%         that develops the topic. Each section is bounded to ~1 page.
% Defer:  derivations, proofs, algorithmic details and numerical examples to
%         Appendices A (Markov), B (Probabilistic Metrics), C (Process Mining
%         details) and D (end-to-end walkthrough).
% =============================================================================

\chapter{Theoretical Foundation}
\label{ch:theoretical-foundation}

This chapter states the formal vocabulary on which PATHCAST is built. Six
concepts are introduced in succession: process mining
(Section~\ref{sec:process-mining}), mining of software repositories
(Section~\ref{sec:msr}), Markov chains (Section~\ref{sec:markov}), Monte
Carlo simulation and the probabilistic forecasting metrics it requires
(Section~\ref{sec:monte-carlo}), machine learning for software engineering
(Section~\ref{sec:ml-se}), and the entropy of trace variants
(Section~\ref{sec:variant-entropy-cap2}). Each section is intentionally short:
its purpose is to fix notation and quote the central property used elsewhere in
the thesis. Derivations, algorithmic internals and worked examples are
deferred to Appendices~A--D. The literature gap that motivates the method is
synthesised in Chapter~\ref{ch:slr}, not here.


% =============================================================================
% 2.1 PROCESS MINING
% =============================================================================
\section{Process Mining}
\label{sec:process-mining}

Process mining extracts behavioural knowledge from event logs recorded by
information systems~\cite{aalst2016process}. Three classes of analysis are
distinguished: \emph{discovery} extracts a model $M$ from a log
$\mathcal{L}$ alone; \emph{conformance} compares $\mathcal{L}$ against a
given $M$; \emph{enhancement} enriches $M$ with information derived from
$\mathcal{L}$. PATHCAST uses all three.

\begin{definition}[Event, Trace, Event Log]
\label{def:event}
\label{def:event-log}
An event $e = (c, a, t, d)$ is a tuple with case identifier
$c \in \mathcal{C}$, activity $a \in \mathcal{A}$, timestamp
$t \in \mathbb{R}^+$ and attribute map $d \in \mathcal{D}$. Projections
$\kappa(e), \alpha(e), \tau(e)$ return case, activity and timestamp. A
trace $\sigma_c = \langle e_1, \ldots, e_n \rangle$ is the
timestamp-ordered sequence of events of case $c$. An event log
$\mathcal{L} = \{\sigma_{c_1}, \ldots, \sigma_{c_N}\}$ is a multiset of
traces. Since each case corresponds to exactly one trace, we use the case
identifier $c$ and its trace $\sigma_c$ interchangeably where no ambiguity
arises.
\end{definition}

\begin{definition}[Alignment and Alignment-Based Fitness]
\label{def:alignment}
\label{def:fitness}
An alignment $\gamma$ between trace $\sigma$ and model $M$ is a sequence
of synchronous, log and model moves~\cite{adriansyah2011alignments}. With
$\gamma^*$ the cost-optimal alignment,
\begin{align}
\text{fitness}(\sigma, M) &= 1 - \frac{\text{cost}(\gamma^*)}
{\text{cost}(\gamma_{\text{worst}})}, \label{eq:fitness} \\
\text{conf}(c) &= \frac{|\{(s_i,m_i)\in\gamma^*_c : s_i=m_i\}|}{|\gamma^*_c|}. \label{eq:case-conformance-cap2}
\end{align}
Log-level fitness is the trace-weighted average. The case-level
conformance $\text{conf}(c)$ is reused as an ML feature in
Chapter~\ref{ch:method}.
\end{definition}

PATHCAST adopts the Inductive Miner with infrequent-behaviour filtering
(IMf)~\cite{leemans2013discovering} because it is among the few discovery
algorithms that guarantee soundness by construction, producing a
block-structured process tree compatible with SDLC workflows. Discovery
quality is quantified along four dimensions---fitness, precision,
generalisation and simplicity~\cite{aalst2016process}; precision and
generalisation are estimated by the ETC metric and state-visit frequency
analysis, respectively. The internals of IMf, the alignment search and
the precision/generalisation estimators are presented in Appendix~A.


% =============================================================================
% 2.2 MINING SOFTWARE REPOSITORIES
% =============================================================================
\section{Mining Software Repositories}
\label{sec:msr}

Mining Software Repositories (MSR) extracts evidence from version control,
issue trackers, code review platforms and CI/CD pipelines~\cite{hassan2008road,
kagdi2007}. The bridge with process mining was established by Cook and
Wolf~\cite{cook1998} and formalised by Poncin et al.~\cite{poncin2011process}:
software events are reframed as sequences belonging to process instances.

Modern toolchains emit heterogeneous, partly overlapping streams:
issue-tracker status transitions (Jira, GitHub Issues), version-control
commits (Git, SVN), code-review activities (PRs, Gerrit) and CI/CD
outcomes (Jenkins, GitHub Actions). The empirical challenge is
integration~\cite{kochhar2019moving, Dixit2021}: a unified event log
requires identifier mapping, temporal alignment and semantic
harmonisation across tools. Curated datasets and collectors---SEART GitHub
Search~\cite{seart2021}, GH~Archive~\cite{gharchive2023},
SmartSHARK~\cite{trautsch2020smartshark},
GrimoireLab~\cite{dueenas2021grimoirelab} and the Public Jira
Dataset~\cite{montgomery2022jira}---reduce the effort required to identify
and collect suitable repositories for analysis, without eliminating known
validity threats (open-source bias, bot noise, incompleteness). The
construction sub-tasks (extraction, case correlation, activity mapping) and
their archaeological caveats (stratigraphy, contextual interpretation,
taphonomic bias~\cite{Zimmermann2005MiningVersionHistories}---i.e.,
distortions caused by the incomplete preservation of historical repository
data) are operationalised by
PATHCAST Stage~1 (Chapter~\ref{ch:method}). Software-specific construction
patterns are catalogued in Appendix~C, Section~C.4.


% =============================================================================
% 2.3 STOCHASTIC PROCESS MODELING WITH MARKOV CHAINS
% =============================================================================
\section{Stochastic Process Modeling with Markov Chains}
\label{sec:markov}

Software delivery evolves through discrete states (Backlog, In~Progress,
Review, Testing, Done) under queuing, rework and variable service times.
Discrete-time Markov chains (DTMCs) provide the mathematical formalism for
the state-transition layer of PATHCAST~\cite{Spedicato2017, Haggstrom2002}.

\begin{definition}[Discrete-Time Markov Chain]
\label{def:dtmc}
A stochastic process $\{X_t\}_{t \in \mathbb{N}}$ on finite state space
$S = \{1, \ldots, m\}$ is a DTMC if
\begin{equation}
\label{eq:markov-property}
P(X_{t+1} = j \mid X_t, \ldots, X_0) = P(X_{t+1} = j \mid X_t).
\end{equation}
The (time-homogeneous) chain is fully described by a row-stochastic
transition matrix $P = [p_{ij}]$, with $p_{ij} = P(X_{t+1}=j \mid X_t=i)$
satisfying $p_{ij} \geq 0$ and $\sum_j p_{ij} = 1$
(Definition~\ref{def:transition-matrix}).
\end{definition}

\begin{definition}[Transition Matrix]
\label{def:transition-matrix}
$P$ as in Definition~\ref{def:dtmc}; multi-step probabilities follow the
Chapman--Kolmogorov relation $P^{(n)} = P^n$.
\end{definition}

\begin{definition}[Absorbing Chain and Fundamental Matrix]
\label{def:absorbing}
For an absorbing chain with $t$ transient and $r$ absorbing states the
canonical form is
\begin{align}
P &= \begin{pmatrix} Q & R \\ \mathbf{0} & I_r \end{pmatrix}, \label{eq:absorbing-form} \\
N &= (I_t - Q)^{-1} = \sum_{k=0}^{\infty} Q^k, \label{eq:fundamental}
\end{align}
where $N_{ij}$ is the expected number of visits to transient state $j$
starting from transient state $i$ before absorption
(Definition~\ref{def:fundamental-matrix}).
\end{definition}

\begin{definition}[Fundamental Matrix]
\label{def:fundamental-matrix}
The matrix $N$ in Equation~\ref{eq:fundamental} is well-defined because
$\rho(Q) < 1$.
\end{definition}

\begin{theorem}[Properties of Absorbing Chains]
\label{thm:absorption}
With $N = (I-Q)^{-1}$:
\emph{(i)} $E[T_i] = (N\mathbf{1})_i$;
\emph{(ii)} $\text{Var}(T_i) = (N(2N_{\text{dg}}-I)\mathbf{1})_i - E[T_i]^2$;
\emph{(iii)} $B = NR$ is the absorption probability matrix. Proofs in
Appendix~A.
\end{theorem}

Counts $c_{ij}$ of observed $i \to j$ transitions yield the maximum
likelihood estimator
\begin{equation}
\label{eq:mle-transition}
\hat{p}_{ij} = \frac{c_{ij}}{\sum_{k} c_{ik}};
\end{equation}
Bayesian smoothing under a symmetric Dirichlet prior is derived in
Appendix~A. Augmenting a DTMC with state-dependent sojourn distributions
$\{H_{ij}\}$ yields a semi-Markov
process~\cite{RoggeSolti2015NonMarkovianSPN}:

\begin{definition}[Semi-Markov Process]
\label{def:semi-markov}
The pair (DTMC $\{X_n\}$, sojourn distributions $\{H_{ij}\}$) under which
the process remains in $X_n$ for a random time drawn from
$H_{X_n,X_{n+1}}$ before transitioning.
\end{definition}

The first-order assumption hides history dependence and time-varying
transitions; PATHCAST partially mitigates both through optional state-space
expansion and ML-based residual correction (Section~\ref{sec:ml-se}),
neither of which fully removes violations of the Markov property.


% =============================================================================
% 2.4 MONTE CARLO SIMULATION AND PROBABILISTIC FORECASTING METRICS
% =============================================================================
\section{Monte Carlo Simulation and Probabilistic Forecasting Metrics}
\label{sec:monte-carlo}

Monte Carlo simulation estimates statistical properties of intractable
systems through repeated random sampling~\cite{metropolis1949monte,
robert2004monte}. In PATHCAST, MC generates synthetic process trajectories
by alternating sojourn draws from $\hat{F}_s$ and next-state draws from
$P[s,:]$ until absorption (Chapter~\ref{ch:method},
Algorithm~\ref{alg:mc-full}).

Convergence rests on two classical results. The strong law of large
numbers gives, for i.i.d.\ $X_1, X_2, \ldots$ with $E[X]=\mu$,
\begin{equation}
\label{eq:lln}
\bar{X}_M = \tfrac{1}{M}\sum_{m=1}^{M} X_m \xrightarrow{a.s.} \mu.
\end{equation}
The central limit theorem yields, with $\text{Var}(X)=\sigma^2 < \infty$,
\begin{equation}
\label{eq:clt}
\frac{\bar{X}_M - \mu}{\sigma/\sqrt{M}} \xrightarrow{d} \mathcal{N}(0,1),
\end{equation}
so for any trajectory functional $g$ the MC estimator and confidence
interval are
\begin{align}
\hat{\mu}_g &= \tfrac{1}{M}\sum_{m=1}^{M} g(\omega^{(m)}), \label{eq:mc-estimator} \\
\text{CI} &= \hat{\mu}_g \pm z_{\alpha/2} \cdot \tfrac{\hat{\sigma}_g}{\sqrt{M}}. \label{eq:mc-ci}
\end{align}
The Monte Carlo convergence rate remains $O(1/\sqrt{M})$ regardless of
dimensionality (although the estimator variance $\sigma^2$ may itself grow
with dimension), and seed control via the Mersenne Twister
makes the simulation reproducible under fixed seeds and implementation
settings.

\paragraph{Forecasting metrics.} A distributional forecast must be
evaluated jointly for sharpness and calibration. The continuous ranked
probability score~\cite{gneiting2007strictly} is the operational
choice in this thesis:

\begin{definition}[CRPS]
\label{def:crps}
For predicted distribution $\hat{F}$ and observed value $y$,
\begin{equation}
\label{eq:crps}
\text{CRPS}(\hat{F}, y) = \int_{-\infty}^{\infty}
   [\hat{F}(t) - \mathbf{1}(t \geq y)]^2 \, dt.
\end{equation}
CRPS is a strictly proper scoring rule~\cite{gneiting2007strictly} and
reduces to mean absolute error when $\hat{F}$ degenerates to a point mass,
so distributional forecasters and point baselines remain comparable on a
common scale. Calibration is summarised by the average calibration error
(ACE)
\begin{equation}
\label{eq:ace}
\text{ACE} = \tfrac{1}{|\mathcal{A}|}
   \sum_{\alpha \in \mathcal{A}}|\hat{p}_\alpha - \alpha|,
\end{equation}
averaged across the quantile checkpoints $\mathcal{A}$. ACE is the
unweighted mean of the absolute deviations between nominal and empirical
coverage at a fixed set of probability levels; unlike the
\emph{expected} calibration error (ECE), which weights bins by their
population, ACE assigns equal weight to each checkpoint and is therefore
more sensitive to miscalibration in sparsely populated tail quantiles---the
regime most relevant to the percentile-based forecasts produced by
PATHCAST. Worked CRPS-versus-MAE derivations
and numerical examples are in Appendix~B.
\end{definition}


% =============================================================================
% 2.5 MACHINE LEARNING FOR SOFTWARE ENGINEERING
% =============================================================================
\section{Machine Learning for Software Engineering}
\label{sec:ml-se}

Machine learning is widely applied to defect prediction, effort
estimation and code quality assessment~\cite{Wahono2015, pachouly2022};
process-oriented prediction (throughput, workflow optimisation) remains
comparatively underexplored~\cite{kotti2023}. PATHCAST formulates outcome
prediction as supervised learning: a binary classifier
$\hat{f}\colon \mathbb{R}^d \to \{0,1\}$ for on-time versus delayed cases,
and a regressor $\hat{f}_T\colon \mathbb{R}^d \to \mathbb{R}^+$ for lead
time, both trained on a feature vector $\mathbf{x}_c \in \mathbb{R}^d$
derived per case. The candidate algorithms (Random
Forest~\cite{breiman2001random}, XGBoost~\cite{chen2016xgboost},
logistic/linear regression baselines, and an LSTM~\cite{tax2017} as a
consolidated predictive-process-monitoring baseline---included to quantify
the gain of process-oriented modelling rather than to compete with more
recent Transformer-based architectures---for the
sequence comparator) and the training protocol (temporal hold-out,
hyperparameter search, SHAP interpretation) are specified in
Chapter~\ref{ch:methodology}. The working hypothesis is that
process-derived features---conformance, rework count, expected remaining
transitions, completion probability, path entropy---add predictive value
beyond traditional MSR features; the feature catalogue and the
PM-only/MSR-only/Combined comparison protocol are in Chapter~\ref{ch:method},
Table~\ref{tab:features}.


% =============================================================================
% 2.6 VARIANT ENTROPY
% =============================================================================
\section{Variant Entropy}
\label{sec:variant-entropy-cap2}

The trace-variant distribution $V = \{v_1, \ldots, v_k\}$ with empirical
frequencies $p(v_i)$ summarises the diversity of execution paths in an
event log. The Shannon entropy
\begin{equation}
\label{eq:variant-entropy-cap2}
H(V) = - \sum_{v \in V} p(v)\, \log_2 p(v)
\end{equation}
quantifies path uncertainty in bits; its normalised form
$H_{\text{norm}}(V) = H(V) / \log_2 |V| \in [0,1]$ is a scale-independent
process-complexity proxy. Higher $H_{\text{norm}}$ signals greater
behavioural heterogeneity and is expected to be associated with lower
forecast confidence; PATHCAST uses it as both a diagnostic indicator at
Stage~2 (Section~\ref{sec:variant-entropy}) and a feature of the ML
component.


% =============================================================================
% 2.7 CHAPTER SYNTHESIS  (one paragraph)
% =============================================================================
\section{Chapter Synthesis}
\label{sec:ch2-synthesis}

The six concepts compose the PATHCAST pipeline as follows: process mining
yields the discovered backbone and the per-case conformance feature;
mining of software repositories provides the heterogeneous event sources
from which that log is built; the Markov chain layer turns the backbone
into a probabilistic state-transition model with closed-form expressions
for expected completion time and absorption probabilities; Monte Carlo
simulation, paired with CRPS and ACE, produces and evaluates the
distributional forecasts; machine learning consumes process-derived
features alongside MSR features to refine case-level outcome predictions;
and the entropy of variants provides a proxy for process complexity and
expected forecasting difficulty. The empirical evidence that this
integration has not, to the best of our knowledge, been previously proposed
for software development is the subject of Chapter~\ref{ch:slr}.
